\subsection*{Procedimiento y operaciones del simulacro E}
La corrección usa los cocientes $\theta=T/T_{ref}$ y $\delta=p/p_{ref}$: $N/\sqrt\theta$ y $\dot m\sqrt\theta/\delta$. El acoplamiento iguala la potencia del compresor con la de la turbina multiplicada por el rendimiento mecánico, y la razón de expansión sale de la caída de temperatura con el rendimiento total-total.

\begin{enumerate}

\item Temperatura total de entrada en el ensayo.

\[\fitmath{\mathrm{Tin}=250\approx 250\,\mathrm{K}}\]

\item Presión total de entrada en el ensayo.

\[\fitmath{\mathrm{pin}=70000\approx 70000\,\mathrm{Pa}}\]

\item Caudal másico medido.

\[\fitmath{m=14\approx 14\,\mathrm{kg/s}}\]

\item Régimen medido.

\[\fitmath{N=10200\approx 10200\,\mathrm{min^{-1}}}\]

\item Cociente de temperatura respecto a la referencia.

\[\fitmath{\mathrm{th}=\frac{\mathrm{Tin}}{288.15}}\]

\[\fitmath{\mathrm{th}=\frac{(250)}{288.15}\approx 0.8676037\,}\]

\item Cociente de presión respecto a la referencia.

\[\fitmath{\mathrm{de}=\frac{\mathrm{pin}}{101325}}\]

\[\fitmath{\mathrm{de}=\frac{(70000)}{101325}\approx 0.6908463\,}\]

\item Régimen corregido.

\[\fitmath{\mathrm{Nc}=\frac{N}{\sqrt{\mathrm{th}}}}\]

\[\fitmath{\mathrm{Nc}=\frac{(10200)}{\sqrt{(0.8676036786)}}\approx 10950.64\,\mathrm{min^{-1}}}\]

\item Caudal corregido.

\[\fitmath{\mathrm{mc}=\frac{m\cdot \sqrt{\mathrm{th}}}{\mathrm{de}}}\]

\[\fitmath{\mathrm{mc}=\frac{(14)\cdot \sqrt{(0.8676036786)}}{(0.6908462867)}\approx 18.87588\,\mathrm{kg/s}}\]

\item Cociente de temperatura del punto equivalente.

\[\fitmath{\mathrm{th}_{2}=\frac{320}{288.15}\approx 1.110533\,}\]

\item Régimen físico con el mismo régimen corregido.

\[\fitmath{N_{2}=\mathrm{Nc}\cdot \sqrt{\mathrm{th}_{2}}}\]

\[\fitmath{N_{2}=(10950.63943)\cdot \sqrt{(1.110532709)}\approx 11539.98\,\mathrm{min^{-1}}}\]

\item Caudal físico a 101325 Pa con el mismo caudal corregido.

\[\fitmath{m_{2}=\frac{\mathrm{mc}}{\sqrt{\mathrm{th}_{2}}}}\]

\[\fitmath{m_{2}=\frac{(18.87588403)}{\sqrt{(1.110532709)}}\approx 17.9119\,\mathrm{kg/s}}\]

\item Caudal del compresor.

\[\fitmath{\mathrm{mcomp}=22\approx 22\,\mathrm{kg/s}}\]

\item Aumento de temperatura total del compresor.

\[\fitmath{\mathrm{dTc}=240\approx 240\,\mathrm{K}}\]

\item Potencia absorbida por el compresor.

\[\fitmath{\mathrm{Pc}=\mathrm{mcomp}\cdot 1005\cdot \mathrm{dTc}}\]

\[\fitmath{\mathrm{Pc}=(22)\cdot 1005\cdot (240)\approx 5306400\,\mathrm{W}}\]

\item Caudal de la turbina con combustible y sangrado.

\[\fitmath{\mathrm{mt}=\left(1+0.02-0.05\right)\cdot \mathrm{mcomp}}\]

\[\fitmath{\mathrm{mt}=\left(1+0.02-0.05\right)\cdot (22)\approx 21.34\,\mathrm{kg/s}}\]

\item Caída de temperatura total requerida en la turbina.

\[\fitmath{\mathrm{dTt}=\frac{\mathrm{Pc}}{0.99\cdot \mathrm{mt}\cdot 1150}}\]

\[\fitmath{\mathrm{dTt}=\frac{(5306400)}{0.99\cdot (21.34)\cdot 1150}\approx 218.41\,\mathrm{K}}\]

\item Temperatura total de entrada a la turbina.

\[\fitmath{T_{04}=1450\approx 1450\,\mathrm{K}}\]

\item Razón de expansión total-total de la turbina.

\[\fitmath{\mathrm{prt}={\left(1-\frac{\mathrm{dTt}}{0.88\cdot T_{04}}\right)}^{\frac{-1.33}{0.33}}}\]

\[\fitmath{\mathrm{prt}={\left(1-\frac{(218.4100077)}{0.88\cdot (1450)}\right)}^{\frac{-1.33}{0.33}}\approx 2.131101\,}\]

\end{enumerate}
